\BeleskeID{09_3-s020}
% element: 09_3-s020:e05
% slide-id: 09_3-s020
Isti izraz za ukupno kašnjenje sada posmatramo kao funkciju broja stepena:
\[t_p=Nt_{p0}\left(1+\frac f\gamma\right)=Nt_{p0}\left(1+\frac{F^{1/N}}\gamma\right).
\]
Pri diferenciranju se $F$ i $\gamma$ uzimaju kao zadati. Za račun izvoda korisno je zapisati $F^{1/N}=\exp(\ln F/N)$, odakle
\[\frac{dF^{1/N}}{dN}=-\frac{F^{1/N}\ln F}{N^2}.
\]
Uslov nultog izvoda daje prikazanu jednačinu. Sa $f=F^{1/N}$ važi $\ln f=\ln F/N$, pa uslov prelazi u $\gamma+f-f\ln f=0$. Numerički, za $\gamma\approx1$, dobija se $f\approx3.6$.
